\documentclass[11pt,a4paper]{article}
\usepackage[T1]{fontenc}
\usepackage{lmodern}
\usepackage[margin=24mm]{geometry}
\usepackage{amsmath,amssymb,mathtools,bm}
\usepackage{booktabs,longtable,array,seqsplit}
\usepackage{xcolor,hyperref}
\hypersetup{colorlinks=true,linkcolor=blue!50!black,citecolor=green!35!black,
 urlcolor=blue!50!black,pdftitle={TimeDependentLattice: Lean blueprint for the paper}}
\setlength{\parindent}{0pt}
\setlength{\parskip}{0.4em}
\setlength{\emergencystretch}{3em}
\renewcommand{\arraystretch}{1.18}
\newcolumntype{L}[1]{>{\raggedright\arraybackslash}p{#1}}
\newcommand{\R}{\mathbb R}
\newcommand{\E}{\mathbb E}
\newcommand{\norm}[1]{\left\lVert#1\right\rVert}
\newcommand{\lean}[1]{\texttt{\expandafter\seqsplit\expandafter{\detokenize{#1}}}}
\newcommand{\HS}{\mathrm{HS}}
\newcommand{\dd}{\,\mathrm d}
\title{\textbf{Lean Formalization Blueprint}\\
 \large Time-dependent lattice rules for SDEs with constant diffusion}
\author{Du Ouyang \and Josef Dick \and Ziyang Ye}
\date{23 October 2026}
\begin{document}
\maketitle
This blueprint accompanies the paper and describes the
current Lean formalization. All theorem, lemma, definition and equation
numbers below are the numbers printed in that paper's PDF.
File paths are relative to \texttt{TimeDependentLattice/}; a file name
locates a declaration and is not an additional namespace component.

\section{Results and input overview}
\label{sec:overview}
The formalization has two branches. The derivative-tree expansion,
summability and Gaussian norm estimates lead to the sampling results.
The SDE and moment inputs lead to Euler strong-error and bias bounds.
The total-error theorems combine the two branches.

\subsection{Main results}
{\small
\begin{longtable}{@{}L{0.23\linewidth}L{0.73\linewidth}@{}}
\toprule Paper result & Final Lean entry\\ \midrule\endhead
Theorem 4.1; (4.2) &
\lean{Sampling/DimensionUniform.lean}:\newline
\lean{theorem4_1_with_paper_constant},
\lean{theorem4_1_exists_pos_uniform}.\\
Theorem 3.1; (3.1)--(3.2) &
\lean{Sampling/PaperTheorems.lean}:\newline
\lean{EulerSamplingInputs.single_real},
\lean{EulerSamplingInputs.median_real}.\\
Corollary 3.4 &
\lean{Sampling/PaperUniformTheorems.lean}:\newline
\lean{EulerSamplingInputs.single_uniform_real},
\lean{EulerSamplingInputs.median_uniform_real}.\\
Theorems 3.7 and 3.8 &
\lean{Bias/EulerTheorems.lean}:\newline
\lean{theorem3_7_coupled_strong}, \lean{theorem3_7_coupled_bias},
\lean{theorem3_8_coupled_strong}, \lean{theorem3_8_coupled_bias}.
The dimension-uniform conclusions are in \lean{Bias/UniformTheorems.lean}.\\
Corollary 3.9; (3.9)--(3.10) &
\lean{Bias/PaperTotalError.lean},
\lean{Bias/PaperMedianTotalError.lean}:\newline
\lean{corollary3_9_single_paper}, \lean{corollary3_9_median_paper},
and their \lean{_half}/\lean{_one} specializations.
The uniform versions are in \lean{Bias/PaperUniformTotalError.lean}.\\
\bottomrule
\end{longtable}}
The final sampling and total-error statements include \lean{MemLp}
membership and bounds for the real-valued \lean{lpNorm}.
Key intermediate lemmas can be located directly in Section~\ref{sec:navigation}.

Section~\ref{sec:inputs} separates the model assumptions, the supplied
objects and mathematical results, and the measurability and integrability
conditions. Each supplied result is stated in
Section~\ref{sec:math-inputs}, together with its use in the paper,
Lean interface and the conclusions that use it.

\section{Model conditions and accepted inputs}
\label{sec:inputs}
\subsection{Objects and model assumptions}
\label{sec:model}
Throughout, we use the paper's convention $T>0$.
\lean{State stateDim} is $\R^{d_x}$ with its Euclidean norm.
A \lean{Diffusion stateDim noiseDim} represents
$\Sigma:\R^{d_w}\to\R^{d_x}$. For the standard basis $(e_r)$, set
\[
 \Theta_p=\sum_{r=1}^{d_w}|\Sigma e_r|_2^p,\qquad
 \sigma=\|\Sigma\|_{\HS},\qquad L=\mathsf B_1=B_0B_1,
 \qquad \mathsf B_2=B_0B_1^2(2!)^\tau.
\]
\lean{FiniteDimension.lean} and \lean{Bias/DiffusionBounds.lean}
identify these quantities and prove (2.9), including
\[
 \Theta_p\le\Theta_{2/3}^{3p/2}\quad(p\ge2/3),\qquad
 \sigma\le M_{2/3}^{3/2}\quad\hbox{if }\Theta_{2/3}\le M_{2/3}.
\]
For $m=2^J$, the structured type \lean{Index J noiseDim} has cardinality
$s=md_w$. The coefficient representation (2.14)--(2.17) defines the
bridge path and the Euler recursion (2.20):
\[
 X^h_{t_{i+1}}(z)=X^h_{t_i}(z)+h\,b(t_i,X^h_{t_i}(z))
       +\Sigma\bigl(W_{t_{i+1}}(z)-W_{t_i}(z)\bigr),
 \quad h=T/m,
\]
with $t_i=ih$, $t_m=T$, and $F_m(z)=g(X_T^h(z))$ as in (2.21).
Throughout, Gaussian integration is written as
$\phi_s(z)\,\dd z$, where $\phi_s(z)=\prod_{j=1}^s\phi(z_j)$.
The symbols $\gamma_u$ denote the function-space subset weights.

The spatial part of Assumption 2.1 is recorded in \lean{GevreyBounds}:
$B_0,B_1,G_0,G_1,\beta_g>0$, $1\le\tau<3/2$, spatial smoothness, and
\[
\begin{aligned}
 \|D_x^n b(t,x)\|_{\rm op}&\le B_0B_1^n(n!)^\tau,
       &&n\ge1,\quad 0\le t\le T,\\
 \|D^n g(x)\|_{\rm op}&\le G_0G_1^n(n!)^\tau e^{\beta_g|x|_2},
       &&n\ge1.
\end{aligned}
\]
The origin bound $|b(t,0)|_2\le K_0$ from (2.3) is supplied separately.
The proved \lean{drift_norm_le_origin_add_lipschitz} gives
$|b(t,x)|_2\le K_0+L|x|_2$.
The paper's drift on $[0,T]$ is represented by a function of real time,
with a spatially smooth extension outside $[0,T]$, for example zero.
The field \lean{drift_smooth} refers to this representative;
the main estimates use only times in $[0,T]$.

The remaining model conditions depend on the conclusion being used:
\begin{itemize}
\item Corollary 3.4 and dimension-uniform bias bounds use Assumptions
2.2--2.3: common model constants, $|x_0|_2\le R_0$ and
$\Theta_{2/3}\le M_{2/3}$.
\item Theorem 3.7 uses the temporal condition (3.4),
$|b(s,x)-b(t,x)|_2\le H_t(1+|x|_2)|s-t|^{1/2}$.
Theorem 3.8 uses (3.7),
$|b(s,x)-b(t,x)|_2\le L_t(1+|x|_2)|s-t|$.
Both conditions are imposed only on $s,t\in[0,T]$.
\item The general total-error result uses Assumption 3.6, equation (3.3):
$|I_m-I|\le C_{\rm bias}m^{-\theta}$ with $\theta>0$.
The half-order and first-order specializations obtain this bound from
Theorems 3.7 and 3.8.
\end{itemize}
The sampling range is prime $N\ge2$, $0<\varepsilon<1/2$;
for the median it is $R=2k-1$, $k\ge1$ and
$1\le q<k/(1-\varepsilon)$. All stochastic inputs below use a
probability space $(\Omega,\mathcal F,\mathbb P)$.

\subsection{Supplied objects and mathematical results}
\label{sec:math-inputs}
Each entry states what is supplied. The \emph{paper location} identifies
where the object or result is used in the paper; the \emph{used to prove}
line names the conclusions obtained from it in Lean.

\paragraph{Sampling objects.}
The sampling interface takes a measurable function
\lean{normalQuantile}$:\R\to\R$ for the standard-normal quantile
transformation, and a measurable nonnegative extended-real function
$a\mapsto e(a)$ for the shift-averaged worst-case error.
Their quantile and Sobolev-space interpretations enter through these
objects and the two error bounds below.

\textit{Paper location:} The quantile transformation in (2.24) and the
worst-case error in (2.27).

\textit{Lean interface:} \lean{normalQuantile}, \lean{quantile_measurable},
\lean{wce} and \lean{wce_measurable} in \lean{EulerSamplingInputs}.

\textit{Used to prove:} Measurability of the lattice estimators and the
sampling and total-error bounds.

\paragraph{Shift-averaged squared-error bound.}
For $B_u$ as in (5.1), the project chooses the auxiliary weights (5.6),
\[
 \gamma_u=\left(\frac{B_u^2}{\rho_\lambda^{|u|}}\right)^{1/(1+\lambda)},
 \qquad u\ne\varnothing,
\]
with the paper's zero-weight convention. The accepted standard bound is
\[
 \E_\Delta|Q_{a,\Delta}(F_m)-I_m|^2
 \le\left(\sum_{u\ne\varnothing}\frac{B_u^2}{\gamma_u}\right)e(a)^2,
\]
for each fixed generating vector $a$.

\textit{Paper location:} In the proof of Proposition 5.3, using
(2.26)--(2.27).

\textit{Lean interface:} \lean{EulerSamplingInputs.shift_bound}.

\textit{Used to prove:} Proposition 5.3, then Theorem 3.1 and
Corollaries 3.4 and 3.9.

\paragraph{Moment bound for the worst-case error.}
For a uniformly chosen generating vector $A$, the cited bound is
\[
 \E_A e(A)^{2\lambda}\le
 \frac1{N-1}\sum_{u\ne\varnothing}\gamma_u^\lambda\rho_\lambda^{|u|}.
\]
The source results are \cite[Theorem~1, Eq.~(13)]{ye2026median} and
\cite[Example~5]{kuo2010}; the constants in (2.28) are fixed in
Section~\ref{sec:sampling-constants}.

\textit{Paper location:} Lemma 2.5, (2.29), used in the proof of
Proposition 5.3.

\textit{Lean interface:} \lean{EulerSamplingInputs.random_wce_bound}.

\textit{Used to prove:} Proposition 5.3, then Theorem 3.1 and
Corollaries 3.4 and 3.9. Proposition 5.3 is proved from the two error
bounds by conditional Jensen's inequality and optimization of the subset
weights, including the zero-weight cases in
\lean{Sampling/Optimization.lean}.

\paragraph{SDE integral equation.}
Let $X$ be the exact process and $\widetilde X_{t_i}$ the coupled Euler
chain driven by the same noise $\Sigma W$. The accepted identities are
\[
 X_0=x_0,\qquad
 X_{t_{i+1}}=X_{t_i}+\int_{t_i}^{t_{i+1}}b(s,X_s)\,\dd s
                  +\Sigma(W_{t_{i+1}}-W_{t_i}),\quad i<m,
\]
almost surely. Existence and uniqueness are cited in the paper from
\cite[Theorem~7.1]{pages2018numerical}.

\textit{Paper location:} Equation (1.1), used in Lemma E.2 and in the
error recurrences in the proofs of Theorems 3.7 and 3.8, including (E.10).

\textit{Lean interface:} \lean{CoupledMomentInputs.solution_initial}
and \lean{CoupledMomentInputs.solution_step}.

\textit{Used to prove:} Theorems 3.7 and 3.8.

\paragraph{State moments from Lemma E.2.}
Lemma E.2 is accepted from the paper, with its explicit constants
(E.2) and (E.4). The consequences used in Lean are
\[
 \forall s\in[0,T],\quad \|X_s\|_{L^2}\le C_X(|x_0|_2,\sigma),
\]
\[
 \E e^{2\beta_g|X_T|_2}
 +\E e^{2\beta_g|\widetilde X_T|_2}
 \le C_{\exp}(2\beta_g,\sigma)e^{2\beta_g e^{LT}|x_0|_2}.
\]
The explicit formulas for $C_X,C_{\exp}$ appear in
Section~\ref{sec:euler-constants}.
Lemma E.1 is not separately formalized: the paper uses it inside the
proof of E.2's exponential estimate.

\textit{Paper location:} Lemma E.2, (E.3) and (E.5), applied in the
proofs of Theorems 3.7 and 3.8 and in the payoff estimate (E.7).

\textit{Lean interface:} \lean{CoupledMomentInputs.exact_moment}
and \lean{CoupledMomentInputs.joint_exponential}.

\textit{Used to prove:} Theorems 3.7 and 3.8, with the paper's
$C_{\rm pay}$ and bias constants.

\paragraph{Increment and remainder integral identities.}
For $t_i\le s\le t_{i+1}$, write
$\delta_i(s)=b(s,X_s)-b(t_i,X_{t_i})$ and
$R_i=\int_{t_i}^{t_{i+1}}\delta_i(s)\,\dd s$.
The accepted $L^2$ identities are
\[
\begin{split}
 X_s-X_{t_i}&=\int_{t_i}^{s}b(r,X_r)\,\dd r+\Sigma(W_s-W_{t_i}),\\
 R_i&=\int_{t_i}^{t_{i+1}}\delta_i(s)\,\dd s.
\end{split}
\]

\textit{Paper location:} In the proof of Theorem 3.7, the increment
estimate following (E.8) and the local remainder estimate before (E.9).

\textit{Lean interface:} \lean{HalfOrderStochasticInputs.increment_eq}
and \lean{HalfOrderStochasticInputs.remainder_eq}.

\textit{Used to prove:} The increment and local remainder bounds,
and Theorem 3.7.

\paragraph{Brownian increment second-moment identity.}
The standard identity, for $t_i\le s\le t_{i+1}$, is
\[
 \E|\Sigma(W_s-W_{t_i})|_2^2=\sigma^2(s-t_i).
\]

\textit{Paper location:} In the proof of Theorem 3.7, the increment
estimate following (E.8).

\textit{Lean interface:} \lean{HalfOrderStochasticInputs.noise_isometry}.

\textit{Used to prove:} Theorem 3.7.

\paragraph{First-order It\^o/Fubini identities.}
For $s\in[t_i,t_{i+1}]$, set
\[
\begin{split}
 V_i(s)&=b(s,X_s)-b(t_i,X_s),\\
 G_i(s)&=D_xb(t_i,X_s)b(s,X_s)
  +\frac12\sum_{r=1}^{d_w}
          D_x^2b(t_i,X_s)[\Sigma e_r,\Sigma e_r].
\end{split}
\]
The supplied decomposition is
\[
 A_i=\int_{t_i}^{t_{i+1}}V_i(s)\,\dd s
      +\int_{t_i}^{t_{i+1}}(t_{i+1}-s)G_i(s)\,\dd s,
 \qquad \sum_{i<j}R_i=\sum_{i<j}A_i+S_j.
\]

\textit{Paper location:} In the proof of Theorem 3.8, the unnumbered
decomposition between (E.10) and (E.11).

\textit{Lean interface:} \lean{OneOrderStochasticInputs.temporal_rep},
\lean{generator_rep}, \lean{A_integral_eq} and \lean{remainder_split},
which identify the random objects and their integrals.

\textit{Used to prove:} Theorem 3.8, through the internally proved
local remainder estimate (E.14).

\paragraph{Martingale property and isometry.}
The process $S_j$ is a square-integrable martingale for the supplied
filtration. Its accepted isometry/orthogonality identity is
\[
 \E|S_j|_2^2=\sum_{i<j}\int_{t_i}^{t_{i+1}}(t_{i+1}-s)^2
        \E\|D_xb(t_i,X_s)\Sigma\|_{\HS}^2\,\dd s.
\]

\textit{Paper location:} In the proof of Theorem 3.8, the variance
calculation immediately before (E.15).

\textit{Lean interface:} \lean{OneOrderStochasticInputs.martingale_property}
and \lean{OneOrderStochasticInputs.martingale_isometry}.

\textit{Used to prove:} The terminal variance bound and Theorem 3.8.

\paragraph{Finite-grid Doob inequality.}
For a square-integrable martingale $S$, the accepted standard inequality is
\[
 \left\|\max_{\ell\le j}|S_\ell|_2\right\|_{L^2}
       \le2\|S_j\|_{L^2}.
\]

\textit{Paper location:} In the proof of Theorem 3.8, the Doob step
yielding (E.15).

\textit{Lean interface:} \lean{FiniteDoobL2}, supplied through
\lean{OneOrderStochasticInputs.doob_l2}.

\textit{Used to prove:} The numerical maximal bound (E.15),
and Theorem 3.8.

\paragraph{Gaussian representation of the Euler terminal state.}
The paper's map is $z\mapsto X_T^h(z)=\widetilde X_{t_m}(\mathcal B_mz)$
with $t_m=T$, as in (2.19). The accepted terminal law states,
for every Borel set $A\subseteq\R^{d_x}$,
\[
 \mathbb P\{\omega:\widetilde X_T(\omega)\in A\}
 =\int_{\R^s}\mathbf1_A\bigl(X_T^h(z)\bigr)\phi_s(z)\,\dd z.
\]

\textit{Paper location:} The representation (2.19)--(2.22), connecting
the coupled Euler expectation to $I_m$ in Corollary 3.9.

\textit{Lean interface:} \lean{hlaw : IdentDistrib} in the half-order
and first-order entries of \lean{Bias/PaperTotalError.lean} and
\lean{Bias/PaperMedianTotalError.lean}.

\textit{Used to prove:} The half-order and first-order specializations
of Corollary 3.9.

\lean{Bias/GaussianExpectation.lean} proves the consequence
\[
 I_m=\int_{\R^s}F_m(z)\phi_s(z)\,\dd z=\E g(\widetilde X_T)
\]
and transports payoff integrability from the coupled Euler variable.

\subsection{Measurability, integrability and representatives}
\label{sec:technical}
Random variables are represented by measurable functions or by elements
of \lean{Lp}; the supplied representatives equal the displayed random
objects almost surely. The common moment input includes
\[
 X_s\in L^2\ (0\le s\le T),\qquad
 \widetilde X_{t_i}\in L^2\ (i\le m),\qquad
 e^{2\beta_g|X_T|_2},\ e^{2\beta_g|\widetilde X_T|_2}\in L^1.
\]
These are \lean{exact_memLp}, \lean{euler_memLp},
\lean{exact_exponential} and \lean{euler_exponential}.
The half-order input also supplies $L^2$ representatives of the drift,
noise increments, state increments, $\delta_i$ and $R_i$; the first-order input
supplies those of $V_i,G_i,A_i,S_j$.

Pathwise drift integrability is
\[
 \int_{t_i}^{t_{i+1}}|b(s,X_s(\omega))|_2\,\dd s<\infty
 \quad\hbox{for almost every }\omega,
\]
recorded by \lean{solution_integrable}.
The Bochner integrability fields are
\[
 \delta_i\in L^1([t_i,t_{i+1}];L^2(\Omega;\R^{d_x}))
 \quad\hbox{(half order)},
\]
\[
 V_i,\quad s\mapsto(t_{i+1}-s)G_i(s)
   \in L^1([t_i,t_{i+1}];L^2(\Omega;\R^{d_x}))
 \quad\hbox{(first order)}.
\]
They are \lean{delta_integrable}, \lean{temporal_integrable} and
\lean{generator_integrable}. In particular, first-order integrability
means strong measurability and
\[
 \int_{t_i}^{t_{i+1}}\|V_i(s)\|_{L^2}\,\dd s<\infty,
 \qquad
 \int_{t_i}^{t_{i+1}}(t_{i+1}-s)\|G_i(s)\|_{L^2}\,\dd s<\infty.
\]
The scalar variance integrand is also integrable:
\[
 s\mapsto(t_{i+1}-s)^2\E\|D_xb(t_i,X_s)\Sigma\|_{\HS}^2
       \in L^1([t_i,t_{i+1}]),
\]
through \lean{stochastic_integrand_integrable}.
These conditions support the integral identities in the proofs of
Theorems 3.7 and 3.8, especially the It\^o/Fubini step before (E.11)
and the variance calculation before (E.15).
The mathematical identities themselves are listed in
Section~\ref{sec:math-inputs}.

\section{Proof guide and paper-to-Lean navigation}
\label{sec:navigation}
The tables locate proofs of paper statements and the main intermediate
steps. A full statement is named directly; a proof step is followed by
its paper location in parentheses. The supplied inputs are listed in
Section~\ref{sec:math-inputs}.

\subsection{Tree expansion and derivative summability}
The mixed derivative uses Mathlib's iterated Fr\'echet derivative on
coordinate directions. Ordered finite partitions are identified with
unordered set partitions; $C^\infty$ smoothness proves invariance under
permutations. The canonical child order therefore has the same value
as the paper's unordered children.
The higher-order chain rule follows the partition formula of
\cite[Theorem~1]{levy2006why}.
The construction of inverse maps between expansion choices and timed
trees gives (B.6). Coefficient counts, time factors and subset sums then
prove Theorem 4.1 for every $\kappa>0$ and $p>2/3$:
\[
 \sup_{m\in\mathcal M}\sum_{u\ne\varnothing}
       \kappa^{|u|}\mathcal D_{m,u}^p
 \le C_{\mathcal D}(\kappa,p;\Theta_{2/3}).
\]
{\small
\begin{longtable}{@{}L{0.33\linewidth}L{0.63\linewidth}@{}}
\toprule Paper result or proof step & Lean proof location\\ \midrule\endhead
Coefficient upper bound and support (in proof of Lemma 2.4; (2.18)) &
\lean{BrownianBridge/Coefficients.lean}:
\lean{coefficient_detail_le}, \lean{coefficient_detail_eq_zero_of_not_mem_supportInterior}.
\lean{BrownianBridge/Path.lean} supplies the linear path.\\
Lemma 4.2, (4.6)--(4.8) &
\lean{Euler/Propagator.lean}:
\lean{compensatedStateCoordinateDerivative_eq_sum_propagator},
\lean{compensatedStateDerivative_eq_sum_nonlinearForcing}.
\lean{Trees/Value.lean}:
\lean{nonlinearStateForcing_eq_sum_driftRootChoices} supplies the local
forcing reindexing used in (4.7).
\lean{Trees/AnalyticBound.lean}: \lean{norm_stepPropagator_le_exp}.\\
Lemma B.1, (B.1) &
\lean{Trees/DerivativeExpansion.lean}:
\lean{paperFinpartitionChainTerm_payoff_eq_sum_directChoices},
\lean{sum_payoff_genericTerm_eq_sum_endpointChoices},
\lean{endpointPartitionDirectEquivPayoffRootChoice}.\\
Definitions B.2--B.4 &
\lean{Trees/PayoffTree.lean}: \lean{LabelledPayoffTree};
\lean{Trees/TimeAssignment.lean};
\lean{Trees/Value.lean}: \lean{TimedPayoffTree.value}.\\
Lemma B.5, (B.6) &
\lean{Trees/DerivativeExpansion.lean}:
\lean{mixedDerivative_integrand_eq_treeExpansionSum}.
The inverse constructions are in \lean{Trees/Correspondence.lean}:
\lean{PayoffExpansionChoice.equivTimedPayoffTree}.\\
Lemma B.6, (B.14)--(B.15) &
\lean{Trees/Counting.lean}: \lean{PayoffTreeData.expression_structure}.\\
Lemma B.7, (B.16) &
\lean{Trees/PlaneCounting.lean}: \lean{weighted_payoffTree_count}.
The Catalan encoding yields $64^k k!$.\\
Lemma B.8, (B.17)--(B.18) &
\lean{Trees/AnalyticBound.lean}:
\lean{TimedPayoffTree.normalized_norm_value_le_factorized},
\lean{analyticCoefficient_sum_le}.\\
Independent-time summation (in proof of Proposition 4.3; (B.20)) &
\lean{Trees/TimeFactorBound.lean}:
\lean{timeAssignment_direction_sum_le_eta_factors}.\\
Proposition 4.3, (4.10) &
\lean{Regularity/Envelope.lean}:
\lean{normalizedDerivativeEnvelope_le_keyMajorantSum}.\\
Lemma 4.4, (4.11)--(4.12) (proof in Appendix C) &
\lean{Regularity/Eta.lean}: \lean{eta_detailIndices_laminar_bound};
\lean{Regularity/DetailSummability.lean}:
\lean{detail_powersetCard_eta_rpow_le}.\\
Lemma 4.5, (4.13)--(4.14) (proof in Appendix C) &
\lean{Regularity/PaperSummability.lean}:
\lean{endpointWeightedSum_le_paperTimeFactorConstant},
\lean{allSubsetTimeFactorSum_le_paperTimeFactorConstant}.\\
$\rho$-summability defining $C_{\rm sum}$ (in proof of Theorem 4.1; (4.17)) &
\lean{Regularity/MajorantSummation.lean}:
\lean{derivativeEnvelopeRhoSum_le_derivativeRhoConstant}.\\
Theorem 4.1 &
\lean{Sampling/DimensionUniform.lean}:
\lean{paperDerivativeConstant_eq_paper_formula},
\lean{paperDerivativeConstant_mono},
\lean{theorem4_1_exists_pos_uniform}.\\
\bottomrule
\end{longtable}}
The $3/2$ laminar-support factor is retained. The constant is independent
of $m$ and $x_0$; common model constants and a bound on $\Theta_{2/3}$
give the dimension-uniform version.

\subsection{Gaussian norm and sampling error}
For $u\ne\varnothing$ and $f=\partial_uF_m$, the code uses the hybrid product law
\[
 \dd\mu_u=\prod_{j\in u}\alpha_0e^{-2\alpha_0|z_j|}\,\dd z_j
             \prod_{j\notin u}\phi(z_j)\,\dd z_j.
\]
The proved conditional-expectation identification gives
\[
 B_u(F_m)=\alpha_0^{-|u|/2}
       \|\E_{\mu_u}[\partial_uF_m\mid z_u]\|_{L^2(\mu_u)}.
\]
Conditioning integrates the inactive Gaussian coordinates. The factor
$\alpha_0^{-|u|}$ in the square cancels the normalized active density,
so this is exactly the iterated-density expression (5.1):
\[
 B_u(F_m)^2=\int_{\R^{|u|}}
    \left|\int_{\R^{s-|u|}}\partial_uF_m(z_u,z_{-u})
                         \phi_{-u}(z_{-u})\,\dd z_{-u}\right|^2
             \prod_{j\in u}e^{-2\alpha_0|z_j|}\,\dd z_u.
\]
The hybrid $L^2$ integrability needed for this identification is proved
for the actual Euler derivative. Conditional-expectation contraction
then gives Lemma 5.2. Weight optimization and the two lattice inputs
produce Proposition 5.3; tail integration for independent copies gives
Lemma 5.4 and the median bound.
{\small
\begin{longtable}{@{}L{0.33\linewidth}L{0.63\linewidth}@{}}
\toprule Paper result or proof step & Lean proof location\\ \midrule\endhead
Gaussian component identification (used in proof of Lemma 5.2; (5.1)) &
\lean{Sampling/GaussianSpacePaper.lean}:
\lean{gaussianSpaceComponent_eq_paper},
\lean{eulerGaussianSpaceComponent_eq_paper}.\\
Lemma 5.1 (proof in Appendix D) &
\lean{Sampling/HybridMomentAssembly.lean}:
\lean{hybridEulerTerminalMoment_le_auxiliaryMomentConstant}.
Both $\Theta_1=0$ and $\Theta_1>0$ are covered.\\
Lemma 5.2, (5.3) &
\lean{Sampling/GaussianSpaceNorm.lean}: \lean{lemma5_2};
\lean{Sampling/EulerTransfer.lean}: \lean{eulerGaussianSpaceComponent_le}.\\
Subset-weight optimization (in proof of Proposition 5.3; (5.6)) &
\lean{Sampling/Optimization.lean}: \lean{auxiliarySubsetWeight}
and the proved balance identities, including zero weights.\\
Proposition 5.3, (5.4)--(5.5) &
\lean{Sampling/Transfer.lean}: \lean{latticeMomentScale};
\lean{Sampling/FractionalMoment.lean}: \lean{proposition5_3_prod}.\\
Lemma 5.4, (5.7) &
\lean{Sampling/MedianTail.lean}: \lean{median_lintegral_rpow_le_paper},
\lean{median_eLpNorm_le_paper}.\\
Uniform single-copy fractional moment (in proof of Theorem 3.1; (5.9)) &
\lean{Sampling/EulerMain.lean}: \lean{euler_lattice_fractional_moment},
combined with \lean{Sampling/PaperConstants.lean}:
\lean{latticeUniformMomentConstant_le_eulerPaperLatticeMomentConstant}.\\
Theorem 3.1; Corollary 3.4 &
\lean{Sampling/PaperTheorems.lean} and
\lean{Sampling/PaperUniformTheorems.lean}; final real-valued entries
are listed in Section~\ref{sec:overview}.\\
\bottomrule
\end{longtable}}
The estimator definitions implement (2.23)--(2.25).
\lean{Sampling/LatticeProbability.lean} constructs the randomization
measures. \lean{Sampling/EulerMedian.lean} proves independence of repeated errors.
The median order condition proves uniqueness and measurability.

\subsection{Euler strong error and payoff bias}
The half-order branch uses the increment equation and (3.4) to bound
$R_i$. The first-order branch estimates the actual Jacobian and Hessian
terms in $A_i$ and the martingale term $S_j$. In particular, the code
proves (E.11)--(E.12), then the terminal variance and (E.15):
\[
 \E|S_m|_2^2\le\frac{L^2\sigma^2T}{3}h^2,\qquad
 \left\|\max_{j\le m}|S_j|_2\right\|_{L^2}
            \le2L\sigma\sqrt{T/3}\,h.
\]
Telescoping the actual SDE/Euler difference and applying discrete
Gronwall yields the strong estimates. The payoff derivative bound and
Cauchy--Schwarz prove (E.7), with the exact $C_{\rm pay}$ of (E.6).
{\small
\begin{longtable}{@{}L{0.33\linewidth}L{0.63\linewidth}@{}}
\toprule Paper result or proof step & Lean proof location\\ \midrule\endhead
Payoff transfer (used in proofs of Theorems 3.7 and 3.8; (E.7)) &
\lean{Bias/PayoffTransfer.lean}:
\lean{payoff_sub_le_endpoint_exponentials},
\lean{payoff_integral_sub_le_payoffDifferenceConstant}.\\
State-increment estimate (in proof of Theorem 3.7; following (E.8)) &
\lean{Bias/LocalEstimates.lean}: \lean{state_increment_lpNorm_of_integral_eq}.\\
Half-order local remainder (in proof of Theorem 3.7; between (E.8) and (E.9)) &
\lean{Bias/EulerTheorems.lean}: \lean{HalfOrderStochasticInputs.local_bound}.\\
Vector error recurrence (in proof of Theorem 3.8; (E.10)) &
\lean{Bias/EulerProcess.lean}: \lean{coupledEuler_error_step},
\lean{coupledEuler_error_eq_sum_ae}.\\
Trace and Hilbert--Schmidt estimates (in proof of Theorem 3.8; (E.11)--(E.12)) &
\lean{Bias/DiffusionBounds.lean}: \lean{drift_trace_norm_le},
\lean{drift_diffusion_hilbertSchmidt_sq_le}.\\
Drift remainder estimate (in proof of Theorem 3.8; (E.14)) &
\lean{Bias/EulerTheorems.lean}: \lean{OneOrderStochasticInputs.local_bound}.\\
Terminal variance and martingale maximal bound (in proof of Theorem 3.8;
display before (E.15) and (E.15)) &
\lean{Bias/EulerTheorems.lean}:
\lean{OneOrderStochasticInputs.martingale_second_moment},
\lean{OneOrderStochasticInputs.martingale_maximal}.\\
Gronwall and strong-error estimate (in proof of Theorem 3.8; (E.17)) &
\lean{Bias/ErrorRecursion.lean}:
\lean{gridMaxNorm_recurrence_of_local_bounds};
\lean{Bias/EulerTheorems.lean}: \lean{theorem3_8_coupled_strong}.\\
Theorems 3.7 and 3.8 &
\lean{Bias/EulerTheorems.lean}: the strong and bias entries in
Section~\ref{sec:overview}.\\
Theorems 3.7(ii) and 3.8(ii) &
\lean{Bias/UniformTheorems.lean}:
\lean{theorem3_7_coupled_strong_uniform},
\lean{theorem3_7_coupled_bias_uniform},
\lean{theorem3_8_coupled_bias_uniform}.\\
\bottomrule
\end{longtable}}
The resulting biases have orders $m^{-1/2}$ and $m^{-1}$.
The core Euler results allow every integer $m>0$; the sampling and total
error entries specialize to the paper's dyadic $m=2^J$.

\subsection{Total error and dimension-uniform constants}
Corollary 3.9 follows by the triangle inequality from the proved sampling
bound and the bias premise (3.3):
\[
\begin{split}
 \|Q_{A,\Delta}^{N,m}-I\|_{L^1}
 &\le C_\varepsilon N^{-1+\varepsilon}+C_{\rm bias}m^{-\theta},\\
 \|M_{R,N,m}-I\|_{L^q}
 &\le C_{\varepsilon,q,R}N^{-1+\varepsilon}+C_{\rm bias}m^{-\theta}.
\end{split}
\]
The \lean{_half}/\lean{_one} entries in
\lean{Bias/PaperTotalError.lean} and
\lean{Bias/PaperMedianTotalError.lean} obtain the bias from Theorems
3.7/3.8 and use the terminal law in Section~\ref{sec:math-inputs} to
identify its Euler expectation with $I_m$.
\lean{Bias/PaperUniformTotalError.lean} uses Corollary 3.4 for the
sampling part. For dimension-uniform total-error constants,
$\theta$ and $C_{\rm bias}$ must also be chosen commonly across
$d_x,d_w$, as required in the final statement of Corollary 3.9.
The uniform Euler theorems supply these choices for $\theta=1/2$ and
$\theta=1$, with a common $H_t$ or $L_t$ respectively.

\section{Constants and parameter dependence}
\label{sec:constants}
This section collects the paper's selected constants. Their positivity,
series convergence and stated monotonicities are proved in the modules
listed below. We abbreviate $x=|x_0|_2$ and $\sigma=\|\Sigma\|_{\HS}$;
all other model parameters are held fixed when monotonicity in $x$ or
$\sigma$ is mentioned.

\subsection{Tree and summability constants}
In Appendix B, $K$ is chosen in the proof of Lemma B.8, and
$C_{\rm maj}$ is chosen after (B.20):
\[
 K=\max\{1,G_0,G_1,e^{LT}B_0,B_1,2^\tau\},\qquad
 C_{\rm maj}=\max\{64K^6\max(1,T),G_1\max(1,G_0)\}.
\]
\lean{Trees/PaperConstants.lean} proves
\lean{GevreyBounds.majorantConstant_eq_paper_formula}; in particular,
$C_{\rm maj}=64K^6\max(1,T)$.

For $2/3<\rho<1/\tau$, $T>0$ and $a,\Theta\ge0$, define the series
chosen at the end of Appendix C:
\[
\begin{split}
 H_*^{\mathcal E}(a;\Theta)
   &=\sum_{n\ge0}(aT^{\rho/2}\Theta)^n(n!)^{\tau\rho-1},\\
 H_*^{\mathcal J}(a;\Theta)
   &=\sum_{n\ge1}(aC_{\rho,T}\Theta)^n n(n!)^{\tau\rho-1}.
\end{split}
\]
Here $C_{\rho,T}$ is the constant of Lemma 4.4 selected in Appendix C,
implemented by \lean{detailSummabilityConstant}.
With $a'=2^{\tau\rho}a$, the paper's constant is
\[
\begin{split}
 C_H(a,\rho,\tau,T;\Theta)
 ={}&H_*^{\mathcal E}(a;\Theta)
 +H_*^{\mathcal J}(a';\Theta)H_*^{\mathcal E}(a';\Theta)\\
 &+T^\rho\bigl(H_*^{\mathcal E}(a;\Theta)-1\bigr).
\end{split}
\]
These are \lean{paperEndpointSeries}, \lean{paperDetailSeries} and
\lean{paperTimeFactorConstant} in \lean{Regularity/PaperSummability.lean}.
They converge, are nondecreasing in $a,\Theta$, and $C_H(a,\rho,\tau,T;0)=1$.
The auxiliary total definitions use $\max(1,C_H)$; the proved equality
on the displayed parameter range removes this maximum.

For the constant in (4.17) and the end of paper Section 4.3, write
\[
\begin{split}
 H&=C_H(\kappa C_{\rm maj}^{\rho},\rho,\tau,T;
                              \vartheta^{3\rho/2}),\\
 C_{\rm sum}(\kappa,\rho;\vartheta)
   &=H\sum_{k\ge0}(C_{\rm maj}^{\rho}H)^k(k!)^{\tau\rho-1},\\
 \rho&=\frac{2/3+\min\{p,1/\tau\}}2,\qquad
 C_{\mathcal D}(\kappa,p;\vartheta)
   =\left[C_{\rm sum}(\kappa^{\rho/p},\rho;\vartheta)\right]^{p/\rho}.
\end{split}
\]
For $\kappa>0$, $p>2/3$ and $\vartheta\ge0$, the final constant is
positive and nondecreasing in $\kappa,\vartheta$.
Its remaining parameters are exactly $\tau,T,B_0,B_1,G_0,G_1$.
\lean{Sampling/DimensionUniform.lean} identifies this formula through
\lean{paperDerivativeConstant_eq_paper_formula} and proves
\lean{paperDerivativeConstant_mono}.

\subsection{Gaussian transfer and sampling constants}
\label{sec:sampling-constants}
Equations (5.2) and (D.3) select
\[
\begin{split}
 \alpha_0&=1+\beta_g e^{LT}\sqrt T\,\Theta_1(2+1/\sqrt2),\qquad
 \lambda_0=2(\alpha_0-1),\\
 K_{\alpha_0}(\lambda_0)&=\max\{2e^{\lambda_0^2/2},\alpha_0\},\\
 C_{\rm aux}(x,\Theta_1)
   &=e^{2\beta_g e^{LT}(x+K_0T)}
       K_{\alpha_0}(\lambda_0)\,2^{(4+3\sqrt2)/(4+\sqrt2)}.
\end{split}
\]
They are defined in \lean{Sampling/HybridMoment.lean};
\lean{Sampling/HybridMomentAssembly.lean} proves Lemma 5.1 with this
constant, including $\Theta_1=0$.
For $0<\varepsilon<1/2$, equations (2.28) and (5.8)--(5.11) give
\[
 \lambda=\frac1{2(1-\varepsilon)},\qquad r_*=1-\frac\varepsilon2,
 \qquad
 C_{\alpha,r}=\frac{\sqrt{2\pi}e^{\alpha^2/(1-r)}}{\pi^{2r}r(1-r)},
\]
\[
 \rho_\lambda=2C_{\alpha_0,r_*}^{\lambda}\zeta(2r_*\lambda),
\]
\[
 C_{\rm lat}(x,\Theta_1,\Theta_{2/3})
 =C_{\rm aux}(x,\Theta_1)^\lambda
 \left[C_{\mathcal D}\!\left(\rho_\lambda^{1/(1+\lambda)},
       \frac{2\lambda}{1+\lambda};\Theta_{2/3}\right)\right]^{1+\lambda},
\]
\[
\begin{split}
 C_\varepsilon&=2^{1-\varepsilon}C_{\rm lat}^{1-\varepsilon},\\
 C_{\varepsilon,q,R}
   &=\left(\frac{2\lambda k}{2\lambda k-q}\right)^{1/q}
       \binom{2k-1}{k}^{1/(2\lambda k)}C_\varepsilon,
 \qquad R=2k-1.
\end{split}
\]
\lean{Sampling/Rate.lean} fixes the worst-case error parameters;
\lean{Sampling/PaperConstants.lean} defines the exact $C_{\rm lat}$;
\lean{Sampling/PaperTheorems.lean} gives the final single and median
constants. The proved comparison (2.9) lifts the intermediate
$\Theta_\rho$ estimate to the displayed $C_{\mathcal D}$.
The actual $\Theta_1$ remains in $\alpha_0$ and $\rho_\lambda$.

With other model parameters fixed, $C_{\rm lat}$ is nondecreasing in
$x,\Theta_1,\Theta_{2/3}$. Corollary 3.4 evaluates the same constant at
\[
 C_{\rm lat}(R_0,M_{2/3}^{3/2},M_{2/3}).
\]
This uses the common model constants of Assumption 2.2 and the
column-sum bound of Assumption 2.3.

\subsection{Euler constants}
\label{sec:euler-constants}
The formulas in \lean{Bias/Constants.lean} are precisely those printed
in Appendix E. Equations (E.2), (E.4), (E.6) and (E.8) give
\[
\begin{split}
 C_X(x,\sigma)&=e^{LT}(x+K_0T+2\sqrt T\,\sigma),\\
 C_{\exp}(v,\sigma)&=10\exp\!\left(ve^{LT}K_0T
                               +v^2e^{2LT}T(1+\sigma^2)\right),\\
 C_{\rm pay}(x,\sigma)&=G_0G_1\sqrt{2C_{\exp}(2\beta_g,\sigma)}
                                      e^{\beta_g e^{LT}x},\\
 C_{\rm inc}(x,\sigma)&=\sqrt T[K_0+LC_X(x,\sigma)]+\sigma.
\end{split}
\]
Equation (E.9) selects the half-order constants
\[
\begin{split}
 C_{{\rm str},1/2}(x,\sigma)
    &=\tfrac23e^{LT}T\sqrt T
        [LC_{\rm inc}(x,\sigma)+H_t(1+C_X(x,\sigma))],\\
 C_{{\rm bias},1/2}(x,\sigma)
    &=C_{\rm pay}(x,\sigma)C_{{\rm str},1/2}(x,\sigma).
\end{split}
\]
Equations (E.13) and (E.16) select
\[
\begin{split}
 C_A(x,\sigma)&=\tfrac12(L_t+LK_0+\tfrac12\mathsf B_2\sigma^2)
                      +\tfrac12(L_t+L^2)C_X(x,\sigma),\\
 C_{\rm str}(x,\sigma)&=e^{LT}[TC_A(x,\sigma)+2L\sigma\sqrt{T/3}],\\
 C_{\rm bias}(x,\sigma)&=TC_{\rm pay}(x,\sigma)C_{\rm str}(x,\sigma).
\end{split}
\]
Their Lean definitions, in the same order, are
\lean{stateMomentConstant}, \lean{stateExponentialPrefactor},
\lean{payoffDifferenceConstant}, \lean{incrementConstant},
\lean{strongHalfConstant}, \lean{biasHalfConstant},
\lean{driftRemainderConstant}, \lean{strongOneConstant} and
\lean{biasOneConstant}.
The strong first-order estimate is $C_{\rm str}h$ as in (E.17);
writing the bias rate in powers of $m^{-1}$ accounts for the factor
$T$ in $C_{\rm bias}$.

These constants are nondecreasing in $x,\sigma$. With common model
constants, a common $H_t$ or $L_t$, $x\le R_0$ and
$\Theta_{2/3}\le M_{2/3}$, the uniform results evaluate the same formulas
at $(R_0,M_{2/3}^{3/2})$.

\begingroup
\small
\bibliographystyle{alpha}
\bibliography{references}
\endgroup
\end{document}
